\uuid{Mtc0}
\titre{Primitives usuelles}
\auteur{hanani ; mzari}
\organisation{exo7}
\licence{CC BY-NC-SA 4.0}
\datecreate{2020-01-20}
\niveau{L1}
\module{Analyse}
\chapitre{Calcul d'intégrales}
\sousChapitre{Autre}
\difficulte{1}
\theme{Calculs d'intégrales, primitives usuelles}
\exo7qcmid{lille-2.integrales.001}

\contenu{
\begin{question}
\qclassification{127.99}
\qtags{motcle=primitives usuelles}
Parmi les égalités suivantes, cocher celles qui sont vraies :
\begin{answers}
\bad{$\displaystyle \int _0^1\frac{\mathrm{d}x}{(x+1)^2}=\frac{1}{(1+1)^2}-\frac{1}{(0+1)^2}=-\frac{3}{4}$.}
\bad{$\displaystyle \int _0^1\frac{\mathrm{d}x}{x+1}=\frac{1}{1+1}-\frac{1}{0+1}=-\frac{1}{2}$.}
\good{$\displaystyle \int _0^1\frac{\mathrm{d}x}{\sqrt{x+1}}=2(\sqrt{2}-1)$.}
\bad{$\displaystyle \int _0^1\sqrt{x+1}\,\mathrm{d}x=\frac{1}{2(\sqrt{2}-1)}$.}
\end{answers}
\vskip2mm
\begin{explanations}
Avec $u=1+x$, on a : $\mathrm{d}u=\mathrm{d}x$, $u(0)=1$, $u(1)=2$ et
$$\int _0^1\frac{\mathrm{d}x}{(x+1)^2}=\int _1^2\frac{\mathrm{d}u}{u^2}=\left[\frac{-1}{u}\right]_1^2=\frac{1}{2}\quad \mbox{ et }\quad \int _0^1\frac{\mathrm{d}x}{x+1}=\int _1^2\frac{\mathrm{d}u}{u}=\Big[\ln u\Big]_1^2=\ln 2.$$
De même, on vérifie que :
$\displaystyle \int _0^1\frac{\mathrm{d}x}{\sqrt{x+1}}=2(\sqrt{2}-1)$ et $\displaystyle \int _0^1\sqrt{x+1}\,\mathrm{d}x=\frac{2}{3}(2\sqrt{2}-1)$.
\end{explanations}
\end{question}
}
